% ============================================================
% 附录A 代码仓库结构与复现说明
% 【负责人：各代码负责人提供各自脚本的运行说明，组长汇总】
% 赛题要求：代码须可复现，并同步提交安装和运行说明。
% ============================================================
\section{代码仓库结构与复现说明}

\subsection{环境与安装}

\begin{lstlisting}[language=bash]
Python 3.12（实测 3.12.14）
pip install -r requirements.txt
# 依赖：numpy 2.3.5 / pandas 3.0.1 / openpyxl 3.1.5
#       python-docx 1.2.0 / pypdf 6.10.0
\end{lstlisting}

\subsection{代码目录与脚本清单}

清洗线脚本已按流水线编号（\texttt{01--11}），全部产物可按序复现：

\begin{table}[H]
\centering
\caption{数据线脚本清单（代码目录）}
\label{tab:scripts}
\small
\begin{tabularx}{\textwidth}{llL}
\toprule
脚本 & 阶段 & 作用 \\
\midrule
01\_数据审计与事件特征.py & 审计 & 数据审计、事件特征与画像统计 \\
02\_逻辑回归基线.py & 基线 & 9 特征逻辑回归五折基线 \\
03\_交付核验与字典.py & 核验 & 交付核验与特征字典生成 \\
04\_2.0样例分析与异常清点.py & 审计 & 2.0 样例分析与异常清点 \\
05\_跨类型关联分析.py & 二阶段 & 跨类型关联分析 \\
06\_生成整合表.py & 派生 & 整合表生成 \\
07/07b/07c/08\_一阶段清洗.py & 一阶段 & 一阶段清洗（含续跑、日历重建与健壮重写版） \\
09\_契约端到端小样本验证.py & 契约 & 契约 v0.1 六台车端到端验证 \\
10\_全量清洗\_契约v0.1.py & 一阶段 & 500 台全量清洗（契约口径） \\
11\_二阶段跨类型标注分级.py & 二阶段 & 跨类型互证标注与分级 \\
\bottomrule
\end{tabularx}
\end{table}

\begin{TODO}
建模线/其余同学脚本请按上表格式追加（脚本名 / 阶段 / 作用 + 一行运行命令）。
\end{TODO}

\subsection{复现步骤}

\begin{TODO}
组长汇总为编号步骤：(1) 数据放置路径约定；(2) 依次运行脚本；
(3) 产物落盘位置；(4) \texttt{submission.csv} 的生成命令；
(5) 与本文档数字对账（行数守恒 3{,}093{,}296、正类 126 等）。
提交前务必在一台干净环境完整跑通一遍并记录耗时。
\end{TODO}
